\documentclass[12pt]{article}
\usepackage[T1]{fontenc}
\usepackage[utf8]{inputenc}
\usepackage{lmodern}
\usepackage{textcomp}
\usepackage{enumitem}
\usepackage{geometry}
\geometry{margin=1in}
\begin{document}
\begin{center}
Answer Key - Sociology - Graduate Level Assessment 14 \
\small (Answers are ordered exactly as in the question paper.)
\end{center}

\section*{Multiple Choice}
\begin{enumerate}[label=\arabic*.]
\item \textbf{Answer:} D
\\ \textit{Options:} A) flexibility was reduced by the introduction of detailed daily worksheets; B) the state had greater control than managers over production processes; C) ownership was transferred from small shareholders to senior managers; D) employment depended on performance rather than trust and commitment
\\ \textit{Note:} Mulholland (1998) posits that privatization shifted the focus of employment relationships from trust and commitment to a performance-based framework, thereby redefining the expectations and accountability between companies and managers.
\item \textbf{Answer:} C
\\ \textit{Options:} A) the content of the media is determined by market forces; B) the subordinate classes are dominated by the ideology of the ruling class; C) the media manipulate 'the masses' as vulnerable, passive consumers; D) audiences make selective interpretations of media messages
\\ \textit{Note:} Mass-society theory posits that the media play a powerful role in shaping public opinion and behavior by treating individuals as passive recipients of information, thereby highlighting their vulnerability to manipulation.
\item \textbf{Answer:} D
\\ \textit{Options:} A) a complex network of interaction at a micro-level; B) a source of conflict, inequality, and alienation; C) an unstable structure of social relations; D) a normative framework of roles and institutions
\\ \textit{Note:} Structural-Functionalists view society as a complex system of interrelated parts, where norms, roles, and institutions work together to maintain social order and stability, thus characterizing it as a normative framework.
\item \textbf{Answer:} C
\\ \textit{Options:} A) shared working conditions in the manufacturing industry; B) the class consciousness of members of the proletariat; C) local communities, extended kinship networks and shared leisure pursuits; D) collective aspirations to move into the middle class
\\ \textit{Note:} The correct answer highlights that traditional working-class identity was shaped by strong local community ties, extensive family connections, and collective recreational activities, which fostered a sense of belonging and social cohesion within the working class.
\item \textbf{Answer:} C
\\ \textit{Options:} A) government 'white paper'; B) confidential medical records; C) household account book; D) the shares register of a business
\\ \textit{Note:} A household account book is a personal document that typically contains private financial information specific to an individual's household, making it a closed-access document as it is not intended for public viewing or scrutiny.
\end{enumerate}

\section*{True / False}
\begin{enumerate}[label=\arabic*.]
\setcounter{enumi}{5}
\item \textbf{Answer:} True
\\ \textit{Options:} A) True; B) False
\\ \textit{Note:} Butler and Stokes (1969) highlighted the significant role of familial and community socialization in shaping political identities and party affiliations, particularly illustrating how working-class individuals are often influenced by the political attitudes and behaviors prevalent in their social environments.
\item \textbf{Answer:} True
\\ \textit{Options:} A) True; B) False
\\ \textit{Note:} The ecological approach to urban sociology emphasizes the spatial distribution of social groups within urban environments, highlighting how these groups interact, compete for resources, and establish their presence in various neighborhoods, thereby supporting the notion of colonization of urban spaces.
\item \textbf{Answer:} False
\\ \textit{Options:} A) True; B) False
\\ \textit{Note:} The print revolution led to the commercialization of newspapers, resulting in private ownership and profit-driven motives rather than public ownership as a shared resource for all citizens.
\item \textbf{Answer:} True
\\ \textit{Options:} A) True; B) False
\\ \textit{Note:} Secularization is characterized by a decline in the influence of religious institutions on societal norms and values, along with a growing skepticism towards traditional religious beliefs, leading to disengagement and disenchantment.
\item \textbf{Answer:} False
\\ \textit{Options:} A) True; B) False
\\ \textit{Note:} The statement is false because, according to various surveys and research studies, cannabis is the most commonly used recreational drug in Britain, significantly outpacing the use of Ecstasy.
\end{enumerate}

\section*{Long Form Response}
\begin{enumerate}[label=\arabic*.]
\setcounter{enumi}{10}
\item \textbf{Answer:} Auguste Comte is often regarded as the father of sociology, as he was the first to systematically study society through a scientific lens, coining the term "sociology" itself. His formulation of positivism emphasized empirical observation and the application of the scientific method to social phenomena, which laid the groundwork for modern social theory. Comte's ideas about social order and progress through stages-specifically, the theological, metaphysical, and positive stages-suggested that societies evolve in a linear fashion, influencing later theorists such as Durkheim and Marx. Furthermore, his belief in the potential for sociology to contribute to social reform underscores the discipline's practical implications, shaping its relevance in addressing contemporary social issues. Thus, Comte's contributions not only established sociology as a distinct field of inquiry but also provided a framework that continues to inform and challenge modern social theories.
\\ \textit{Note:} correct because it accurately identifies Auguste Comte as the founder of sociology and highlights his pivotal contribution of positivism, which advocates for the scientific study of society and has significantly influenced the development of modern social theory.
\item \textbf{Answer:} A population pyramid that is wider at the bottom than at the top indicates a youthful population structure, typically associated with high birth rates and declining death rates over several generations. This demographic pattern suggests that, historically, societies have experienced limited access to family planning and healthcare, leading to larger families. As death rates decrease due to improvements in healthcare and living conditions, the base of the pyramid expands, reflecting a growing number of young individuals. The implications of this structure can include increased pressure on resources, education systems, and employment opportunities, as a larger youth population enters adulthood. Furthermore, such a demographic profile may contribute to economic growth in the short term but poses challenges in terms of sustainability and social support systems as the population ages.
\\ \textit{Note:} correct because a population pyramid that is wider at the bottom than at the top signifies a youthful demographic characterized by high birth rates and declining mortality rates, indicating historical trends of limited family planning, increased healthcare access, and improvements in living conditions.
\end{enumerate}

\vfill
\noindent\textit{End of Answer Key}
\end{document}